\paragraph{Dataset construction and output interpretation.}
The frozen manifest contains 40 search examples per class and 50 held-out examples
per class. Search articles come from the AG News training split and held-out
articles from its test split. The four class names are World, Sports, Business,
and Sci/Tech. Disjointness is checked both by stored sample IDs and by
source-split/index pairs. Stratification controls class counts in these small
samples, but does not remove all variation in article difficulty or topic mix.
The reported accuracy is ordinary micro-accuracy on a class-balanced fixed set;
per-class errors are not reconstructed from the winner summaries.

Each rendered prompt is supplied as a user message through the target model's
chat template, with a generation prompt added. Greedy decoding and an eight-token
output cap are used throughout. The parser examines the first output line,
limited to 80 characters, attempts to extract a label field if present, and
matches a fixed set of label patterns. Its pattern order is Sci/Tech, Sports,
Business, then World; a numeric label from one to four is a fallback. Outputs
that cannot be parsed are counted as incorrect. Measured fitness
reflects the combination of instruction, model response, decoding cap and parser,
rather than abstract classification ability independently of response format.

This interaction is relevant for the format-oriented prompts. A request for JSON
or explanatory text can consume the short output allowance, while the parser
accepts some variants beyond the canonical label strings. Changing the parser or
the decoding limit after seeing results would change the evaluation task. The
current study holds them fixed and does not infer that a longer explanation
would improve or worsen accuracy on this data.
